\documentclass[11pt,a4paper]{article}
\usepackage[T1]{fontenc}
\usepackage[utf8]{inputenc}
\usepackage{lmodern}
\usepackage[margin=25mm]{geometry}
\usepackage{microtype}
\usepackage{amsmath,booktabs,longtable,array}
\usepackage[round,authoryear]{natbib}
\usepackage{xcolor,xurl}
\usepackage[colorlinks=true,linkcolor=blue!45!black,citecolor=blue!45!black,urlcolor=blue!45!black]{hyperref}
\usepackage{enumitem}
\setlist{nosep,leftmargin=*}
\setlength{\emergencystretch}{2em}
\setlength{\bibsep}{4pt}
\setlength{\parskip}{3pt}
\widowpenalty=10000
\clubpenalty=10000
\newcommand{\ILY}{\emph{I love you}}
\newcommand{\snapshot}{ef7ea7933d8aa02b176d8fbddc496e55a41a7144}
\newif\ifincludesupplement
\includesupplementfalse % Set true to append the source and annotation register.
\newcommand{\caseref}[1]{\footnote{\csname case@#1\endcsname}}
\expandafter\def\csname case@nansen\endcsname{Nansen: continuing union: \href{https://github.com/mannyrayner/let_me_count_the_ways/blob/ef7ea7933d8aa02b176d8fbddc496e55a41a7144/results/annotation/canonical_31_v0_11_v0_3_1/annotations/nansen-maria-93ea665425fd352f/d5338f7ac0dbdd0b998461aa6ff8cd9fbc280d1a43db67ca5e57088f6d7d8c75/request.json}{passage 1}, \href{https://github.com/mannyrayner/let_me_count_the_ways/blob/ef7ea7933d8aa02b176d8fbddc496e55a41a7144/results/annotation/canonical_31_v0_11_v0_3_1/annotations/nansen-maria-93ea665425fd352f/d5338f7ac0dbdd0b998461aa6ff8cd9fbc280d1a43db67ca5e57088f6d7d8c75/output.json}{annotation 1}; \href{https://github.com/mannyrayner/let_me_count_the_ways/blob/ef7ea7933d8aa02b176d8fbddc496e55a41a7144/results/annotation/canonical_31_v0_11_v0_3_1/annotations/nansen-maria-a95aa1804441024d/6dbeb2f2a84e2012c634ccc046d6adc3bc5a745d7ade75aed3dd68ee5778a2b0/request.json}{passage 2}, \href{https://github.com/mannyrayner/let_me_count_the_ways/blob/ef7ea7933d8aa02b176d8fbddc496e55a41a7144/results/annotation/canonical_31_v0_11_v0_3_1/annotations/nansen-maria-a95aa1804441024d/6dbeb2f2a84e2012c634ccc046d6adc3bc5a745d7ade75aed3dd68ee5778a2b0/output.json}{annotation 2}.}
\expandafter\def\csname case@jane\endcsname{Jane: sisterly affection: \href{https://github.com/mannyrayner/let_me_count_the_ways/blob/ef7ea7933d8aa02b176d8fbddc496e55a41a7144/results/annotation/canonical_31_v0_11_v0_3_1/annotations/bronte-jane-eyre-5da4cfc02f8bc2c8/c03369f533e2232b1ded4217e8d6f72de5f1facc2fc907eff489feec9fec6e1a/request.json}{passage 1}, \href{https://github.com/mannyrayner/let_me_count_the_ways/blob/ef7ea7933d8aa02b176d8fbddc496e55a41a7144/results/annotation/canonical_31_v0_11_v0_3_1/annotations/bronte-jane-eyre-5da4cfc02f8bc2c8/c03369f533e2232b1ded4217e8d6f72de5f1facc2fc907eff489feec9fec6e1a/output.json}{annotation 1}.}
\expandafter\def\csname case@rhoda\endcsname{Rhoda: withholding assent: \href{https://github.com/mannyrayner/let_me_count_the_ways/blob/ef7ea7933d8aa02b176d8fbddc496e55a41a7144/results/annotation/commitment_extension_5_v0_12_v0_3_1/annotations/gissing-the-odd-women-912b570ed81fd9e9/2e4a6bc03c1777eb92f6d9fb84369fa6cdf6d5a94873511aeb5406bf99b8040f/request.json}{passage 1}, \href{https://github.com/mannyrayner/let_me_count_the_ways/blob/ef7ea7933d8aa02b176d8fbddc496e55a41a7144/results/annotation/commitment_extension_5_v0_12_v0_3_1/annotations/gissing-the-odd-women-912b570ed81fd9e9/2e4a6bc03c1777eb92f6d9fb84369fa6cdf6d5a94873511aeb5406bf99b8040f/output.json}{annotation 1}.}
\expandafter\def\csname case@mary\endcsname{Mary: conditional refusal: \href{https://github.com/mannyrayner/let_me_count_the_ways/blob/ef7ea7933d8aa02b176d8fbddc496e55a41a7144/results/annotation/canonical_31_v0_11_v0_3_1/annotations/eliot-middlemarch-da9a5c10fa10ae23/90b4631af4c322e9416dd4dbc176c9b9b841c9d56498d9e7a4fabcc3f8120775/request.json}{passage 1}, \href{https://github.com/mannyrayner/let_me_count_the_ways/blob/ef7ea7933d8aa02b176d8fbddc496e55a41a7144/results/annotation/canonical_31_v0_11_v0_3_1/annotations/eliot-middlemarch-da9a5c10fa10ae23/90b4631af4c322e9416dd4dbc176c9b9b841c9d56498d9e7a4fabcc3f8120775/output.json}{annotation 1}.}
\expandafter\def\csname case@nora\endcsname{Nora: changed feelings: \href{https://github.com/mannyrayner/let_me_count_the_ways/blob/ef7ea7933d8aa02b176d8fbddc496e55a41a7144/results/annotation/canonical_31_v0_11_v0_3_1/annotations/ibsen-et-dukkehjem-22aa549e55c520e4/d6ccdcad59be3fd7f845ef49fa1b1b5cdbd9f46d8e8d001b4dd544b0aa2be3cd/request.json}{passage 1}, \href{https://github.com/mannyrayner/let_me_count_the_ways/blob/ef7ea7933d8aa02b176d8fbddc496e55a41a7144/results/annotation/canonical_31_v0_11_v0_3_1/annotations/ibsen-et-dukkehjem-22aa549e55c520e4/d6ccdcad59be3fd7f845ef49fa1b1b5cdbd9f46d8e8d001b4dd544b0aa2be3cd/output.json}{annotation 1}.}
\expandafter\def\csname case@lorck\endcsname{Lorck: unwanted offering: \href{https://github.com/mannyrayner/let_me_count_the_ways/blob/ef7ea7933d8aa02b176d8fbddc496e55a41a7144/results/annotation/commitment_extension_5_v0_12_v0_3_1/annotations/skram-constance-ring-2b09d5b293401db7/2f475252a00424783a58e13cc6446a3983579650bcdd4589d6adf88e6b797fc4/request.json}{passage 1}, \href{https://github.com/mannyrayner/let_me_count_the_ways/blob/ef7ea7933d8aa02b176d8fbddc496e55a41a7144/results/annotation/commitment_extension_5_v0_12_v0_3_1/annotations/skram-constance-ring-2b09d5b293401db7/2f475252a00424783a58e13cc6446a3983579650bcdd4589d6adf88e6b797fc4/output.json}{annotation 1}.}
\expandafter\def\csname case@henny\endcsname{Henny: acceptance: \href{https://github.com/mannyrayner/let_me_count_the_ways/blob/ef7ea7933d8aa02b176d8fbddc496e55a41a7144/results/annotation/commitment_extension_5_v0_12_v0_3_1/annotations/skram-lucie-1843e92d8f2ba6e5/268f9d493906123556796da41b2f448c1137d48b559d837a82b6e2896ecb4fb8/request.json}{passage 1}, \href{https://github.com/mannyrayner/let_me_count_the_ways/blob/ef7ea7933d8aa02b176d8fbddc496e55a41a7144/results/annotation/commitment_extension_5_v0_12_v0_3_1/annotations/skram-lucie-1843e92d8f2ba6e5/268f9d493906123556796da41b2f448c1137d48b559d837a82b6e2896ecb4fb8/output.json}{annotation 1}.}
\expandafter\def\csname case@widdowson\endcsname{Widdowson: courtship terms: \href{https://github.com/mannyrayner/let_me_count_the_ways/blob/ef7ea7933d8aa02b176d8fbddc496e55a41a7144/results/annotation/commitment_extension_5_v0_12_v0_3_1/annotations/gissing-the-odd-women-9c7d43c0683c5c32/d7f4ecf135218d9aac9e484baddcdf27e602434ce4368727d00d85fe586b8fa9/request.json}{passage 1}, \href{https://github.com/mannyrayner/let_me_count_the_ways/blob/ef7ea7933d8aa02b176d8fbddc496e55a41a7144/results/annotation/commitment_extension_5_v0_12_v0_3_1/annotations/gissing-the-odd-women-9c7d43c0683c5c32/d7f4ecf135218d9aac9e484baddcdf27e602434ce4368727d00d85fe586b8fa9/output.json}{annotation 1}.}
\expandafter\def\csname case@lorckletter\endcsname{Lorck: considered letter: \href{https://github.com/mannyrayner/let_me_count_the_ways/blob/ef7ea7933d8aa02b176d8fbddc496e55a41a7144/results/annotation/commitment_extension_5_v0_12_v0_3_1/annotations/skram-constance-ring-51a22b8927b39e5a/ef2906b62b9f30ee004c6b1c0de4f199ff10e5817d4dc3f15ccad83d7f1c44b0/request.json}{passage 1}, \href{https://github.com/mannyrayner/let_me_count_the_ways/blob/ef7ea7933d8aa02b176d8fbddc496e55a41a7144/results/annotation/commitment_extension_5_v0_12_v0_3_1/annotations/skram-constance-ring-51a22b8927b39e5a/ef2906b62b9f30ee004c6b1c0de4f199ff10e5817d4dc3f15ccad83d7f1c44b0/output.json}{annotation 1}.}
\expandafter\def\csname case@boarham\endcsname{Boarham: conditional offer: \href{https://github.com/mannyrayner/let_me_count_the_ways/blob/ef7ea7933d8aa02b176d8fbddc496e55a41a7144/results/annotation/canonical_31_v0_11_v0_3_1/annotations/bronte-tenant-of-wildfell-hall-0e20a0f36fe0556a/f4aba1105ded751735e4ce983961902bbfe5e3e0f182daf0a6a483114d72d5fc/request.json}{passage 1}, \href{https://github.com/mannyrayner/let_me_count_the_ways/blob/ef7ea7933d8aa02b176d8fbddc496e55a41a7144/results/annotation/canonical_31_v0_11_v0_3_1/annotations/bronte-tenant-of-wildfell-hall-0e20a0f36fe0556a/f4aba1105ded751735e4ce983961902bbfe5e3e0f182daf0a6a483114d72d5fc/output.json}{annotation 1}.}
\expandafter\def\csname case@emmafuture\endcsname{Emma: prospective inducement: \href{https://github.com/mannyrayner/let_me_count_the_ways/blob/ef7ea7933d8aa02b176d8fbddc496e55a41a7144/results/annotation/canonical_31_v0_11_v0_3_1/annotations/flaubert-madame-bovary-baeec739adc7a033/1287a4d6fc09b57c5bf533b7ac91e3606f3244da9ec34d17e1b44476a01c4bac/request.json}{passage 1}, \href{https://github.com/mannyrayner/let_me_count_the_ways/blob/ef7ea7933d8aa02b176d8fbddc496e55a41a7144/results/annotation/canonical_31_v0_11_v0_3_1/annotations/flaubert-madame-bovary-baeec739adc7a033/1287a4d6fc09b57c5bf533b7ac91e3606f3244da9ec34d17e1b44476a01c4bac/output.json}{annotation 1}.}
\expandafter\def\csname case@conscience\endcsname{Constance: self-persuasion: \href{https://github.com/mannyrayner/let_me_count_the_ways/blob/ef7ea7933d8aa02b176d8fbddc496e55a41a7144/results/annotation/commitment_extension_5_v0_12_v0_3_1/annotations/skram-constance-ring-342b3cad3e79af40/c3dab5fcf6ba007a29018da029714f02cbdacd832d72dc9dbd5090e1fa90bad6/request.json}{passage 1}, \href{https://github.com/mannyrayner/let_me_count_the_ways/blob/ef7ea7933d8aa02b176d8fbddc496e55a41a7144/results/annotation/commitment_extension_5_v0_12_v0_3_1/annotations/skram-constance-ring-342b3cad3e79af40/c3dab5fcf6ba007a29018da029714f02cbdacd832d72dc9dbd5090e1fa90bad6/output.json}{annotation 1}.}
\expandafter\def\csname case@lucy\endcsname{Lucy: private vows: \href{https://github.com/mannyrayner/let_me_count_the_ways/blob/ef7ea7933d8aa02b176d8fbddc496e55a41a7144/results/annotation/commitment_extension_5_v0_12_v0_3_1/annotations/von-arnim-vera-2ec5625a94895d2d/80ae94ece0a69269808a4baa0172b8bb3db889290dab4a2d22dc3c68ad341a64/request.json}{passage 1}, \href{https://github.com/mannyrayner/let_me_count_the_ways/blob/ef7ea7933d8aa02b176d8fbddc496e55a41a7144/results/annotation/commitment_extension_5_v0_12_v0_3_1/annotations/von-arnim-vera-2ec5625a94895d2d/80ae94ece0a69269808a4baa0172b8bb3db889290dab4a2d22dc3c68ad341a64/output.json}{annotation 1}.}
\expandafter\def\csname case@nikki\endcsname{Nikki: words escaping: \href{https://github.com/mannyrayner/let_me_count_the_ways/blob/ef7ea7933d8aa02b176d8fbddc496e55a41a7144/results/annotation/canonical_31_v0_11_v0_3_1/annotations/cofield-nikkis-touch-be393a6a3d4f0015/8748d81cfb9633467c173fee77c0531c152cc38672ab0895b0905c7fcc4efccd/request.json}{passage 1}, \href{https://github.com/mannyrayner/let_me_count_the_ways/blob/ef7ea7933d8aa02b176d8fbddc496e55a41a7144/results/annotation/canonical_31_v0_11_v0_3_1/annotations/cofield-nikkis-touch-be393a6a3d4f0015/8748d81cfb9633467c173fee77c0531c152cc38672ab0895b0905c7fcc4efccd/output.json}{annotation 1}.}
\expandafter\def\csname case@wentworth\endcsname{Wentworth: renewed offer: \href{https://github.com/mannyrayner/let_me_count_the_ways/blob/ef7ea7933d8aa02b176d8fbddc496e55a41a7144/results/annotation/canonical_31_v0_11_v0_3_1/annotations/austen-persuasion-473794649dc5ac56/da01e0aca4075dff2d35f51a87ae5ee7389719960022a5a19e050890f8763c33/request.json}{passage 1}, \href{https://github.com/mannyrayner/let_me_count_the_ways/blob/ef7ea7933d8aa02b176d8fbddc496e55a41a7144/results/annotation/canonical_31_v0_11_v0_3_1/annotations/austen-persuasion-473794649dc5ac56/da01e0aca4075dff2d35f51a87ae5ee7389719960022a5a19e050890f8763c33/output.json}{annotation 1}.}
\expandafter\def\csname case@emmaexpressive\endcsname{Emma: emotional overflow: \href{https://github.com/mannyrayner/let_me_count_the_ways/blob/ef7ea7933d8aa02b176d8fbddc496e55a41a7144/results/annotation/canonical_31_v0_11_v0_3_1/annotations/flaubert-madame-bovary-9c9d14e73da47022/77d181c48ecebf02bc5d8daa7fa51177acb5ea992f5c08cf53a04eb6fb55d7e1/request.json}{passage 1}, \href{https://github.com/mannyrayner/let_me_count_the_ways/blob/ef7ea7933d8aa02b176d8fbddc496e55a41a7144/results/annotation/canonical_31_v0_11_v0_3_1/annotations/flaubert-madame-bovary-9c9d14e73da47022/77d181c48ecebf02bc5d8daa7fa51177acb5ea992f5c08cf53a04eb6fb55d7e1/output.json}{annotation 1}; \href{https://github.com/mannyrayner/let_me_count_the_ways/blob/ef7ea7933d8aa02b176d8fbddc496e55a41a7144/results/annotation/canonical_31_v0_11_v0_3_1/annotations/flaubert-madame-bovary-c648bd2e9a7ad647/fbfa1810e090f58e0094985a122d76b8cda1aa061d117e1a0cddc8476e2bbb99/request.json}{passage 2}, \href{https://github.com/mannyrayner/let_me_count_the_ways/blob/ef7ea7933d8aa02b176d8fbddc496e55a41a7144/results/annotation/canonical_31_v0_11_v0_3_1/annotations/flaubert-madame-bovary-c648bd2e9a7ad647/fbfa1810e090f58e0094985a122d76b8cda1aa061d117e1a0cddc8476e2bbb99/output.json}{annotation 2}.}
\expandafter\def\csname case@stjohn\endcsname{St John: a look verbalized: \href{https://github.com/mannyrayner/let_me_count_the_ways/blob/ef7ea7933d8aa02b176d8fbddc496e55a41a7144/results/annotation/canonical_31_v0_11_v0_3_1/annotations/bronte-jane-eyre-8d223e8c973ddc84/afa503382c40d1751604ad37af8fa33bf0cfd99cc42fbfe2674025d437bd27b2/request.json}{passage 1}, \href{https://github.com/mannyrayner/let_me_count_the_ways/blob/ef7ea7933d8aa02b176d8fbddc496e55a41a7144/results/annotation/canonical_31_v0_11_v0_3_1/annotations/bronte-jane-eyre-8d223e8c973ddc84/afa503382c40d1751604ad37af8fa33bf0cfd99cc42fbfe2674025d437bd27b2/output.json}{annotation 1}.}
\expandafter\def\csname case@birkin\endcsname{Birkin: the formula's limits: \href{https://github.com/mannyrayner/let_me_count_the_ways/blob/ef7ea7933d8aa02b176d8fbddc496e55a41a7144/results/annotation/canonical_31_v0_11_v0_3_1/annotations/lawrence-women-in-love-96fc35925c56519b/2bd6ba5b9158b1b4ede8ee93276693e77a37b8ed38f16482cfb232d807c38bab/request.json}{passage 1}, \href{https://github.com/mannyrayner/let_me_count_the_ways/blob/ef7ea7933d8aa02b176d8fbddc496e55a41a7144/results/annotation/canonical_31_v0_11_v0_3_1/annotations/lawrence-women-in-love-96fc35925c56519b/2bd6ba5b9158b1b4ede8ee93276693e77a37b8ed38f16482cfb232d807c38bab/output.json}{annotation 1}; \href{https://github.com/mannyrayner/let_me_count_the_ways/blob/ef7ea7933d8aa02b176d8fbddc496e55a41a7144/results/annotation/canonical_31_v0_11_v0_3_1/annotations/lawrence-women-in-love-86dea678c8044a4f/c4dfafbe453c344111cfa4a80a4c089ae7bea503bb8690dccd520e6bd201bcbd/request.json}{passage 2}, \href{https://github.com/mannyrayner/let_me_count_the_ways/blob/ef7ea7933d8aa02b176d8fbddc496e55a41a7144/results/annotation/canonical_31_v0_11_v0_3_1/annotations/lawrence-women-in-love-86dea678c8044a4f/c4dfafbe453c344111cfa4a80a4c089ae7bea503bb8690dccd520e6bd201bcbd/output.json}{annotation 2}.}

\hypersetup{pdftitle={Let Me Count the Ways: What Does I Love You Commit Its Speaker To?},pdfauthor={ChatGPT C-LARA-Instance and Manny Rayner}}

\title{Let Me Count the Ways:\\What Does “I Love You” Commit Its Speaker To?}
\author{ChatGPT C-LARA-Instance \and Manny Rayner}
\date{Working draft v0.4 --- 20 September 2026}

\begin{document}
\maketitle

\begin{abstract}
Disagreement about what “I love you” means can become disagreement about what one person owes another. A subsequent loss of love has different implications if the earlier declaration avowed a feeling, made an undertaking, or gave words to an affective impulse. We investigate these possibilities as three nonexclusive dimensions: truth-presenting avowal (T), relational undertaking (P), and expressive/reflexive production (E). An automated analysis of 225 retained occurrences from 30 complete literary works in seven languages is supplemented by a targeted five-novel extension yielding 27 further occurrences. Contextualized readings support distinctions between avowing affection, undertaking a relationship, and words represented as escaping inhibition. The extension also tests an initial expectation: its declarations under pressure do not produce low-T/high-P classifications. This result brings out the difference between avowing love and actually feeling it, and between making an undertaking and sustaining a relationship through reassurance. The contribution is a body of textual evidence and an inspectable procedure for evaluating consequential distinctions, rather than a claim that three categories exhaust love language. The classifications leave open finer questions about the loving relation, the precise undertaking, and disagreements between speakers and hearers.
\end{abstract}

\section{The puzzle: what have the words committed us to?}

Someone says “I love you” and later says, “That was then. Now I don't love you any more.” Has a promise been broken, or have feelings changed? Both parties might agree that the earlier declaration was sincere and still disagree about this. One understands it as having established an obligation that survives a change of feeling; the other understands it as having described how things stood at the time. Their disagreement concerns what the declaration did, as well as whether it was true.

A third possibility complicates the exchange. The speaker might say, “I'm sorry. It just came out.” The claim is that an emotional state produced the words before they could be considered. Whether that explanation is credible, and what it excuses, are further questions. A hearer may nevertheless have understood those words as a commitment and made consequential decisions in reliance on them. The difference can affect expectations of fidelity, care, marriage, or a shared future. It can turn a painful change into what someone experiences as betrayal.

The distinction between an undertaking and the other possibilities is therefore especially consequential. If a promise was made, we must ask what it concerned, what conditions attached to it, and whether it was honoured. If a feeling was avowed, the questions concern sincerity and whatever responsibility people have for changing feelings. If words emerged as an affective response, deliberation and control become relevant. These are different grounds for ethical assessment; none supplies a complete verdict. A promise need not be unconditional, and the absence of a promise does not remove every responsibility towards another person.

The familiar complaint “I don't know what love means” can concern these consequences of speaking as much as the feeling itself. A dictionary definition does not tell a couple whether a particular declaration promised permanence. If speakers encounter the formula in avowals, pledges, reassurances and emotional outbursts, learning from those uses need not give them an agreed rule for a new occasion. This is a possible explanation of the uncertainty, not a developmental finding established by the present study. Competent adults, as well as young learners, can differ over intention, reasonable uptake and retrospective interpretation.

There is an academic puzzle too. In the “I Love You” fragment of \emph{A Lover's Discourse}, Roland Barthes describes the formula as having “no other referent than its utterance: it is a performative” \citep[p.~148]{barthes1978}. Yet many readily imagined uses seem to answer a question about a feeling, while others seem to make a commitment. Barthes's stated focus is the repeated love cry after the first avowal \citep[p.~147]{barthes1978}; his account of performativity also differs from the narrower undertaking category developed below. These qualifications matter. Our question is how far contextual variation warrants a differentiated account, rather than whether an isolated example refutes his entire treatment.

The proposal develops an earlier exploratory essay that assembled examples from novels and films \citep{rayner2014}. Plausible examples, however, are only a beginning. Invented ones may build the intended interpretation into their descriptions. Memorable quotations may be exceptional, or depend on omitted context. Perhaps one interpretation dominates almost completely; perhaps apparently different functions nearly always coincide. Selection and interpretation can both favour the proposed theory.

We address this evidential problem by applying explicit extraction and annotation criteria across complete works. Our initial corpus contains 30 novels and plays in seven languages, yielding 225 retained occurrences. A subsequent, deliberately targeted extension adds five novels and 27 occurrences. We keep these stages distinguishable because the later selection was motivated by a question raised during interpretation. The aim is to establish whether broad distinctions between avowal, undertaking and affective production help explain actual passages, and to make the evidence and procedure accessible to criticism. Establishing that first level would be progress even without a complete account of love or its obligations.

\section{Three questions to ask of a declaration}
\label{sec:dimensions}

We use ILY as a convenient name for the expressions under investigation. T, P and E identify three questions about an occurrence, not three mutually exclusive dictionary senses. Familiar examples help introduce them, but should not be mistaken for independently classified corpus data. The following discussion distinguishes these introductory illustrations from examples with recorded scores.

\paragraph{T: does the utterance avow or present a loving state as true?}
In \emph{Le Petit Prince}, the prince is leaving the small planet where he has cared for a demanding rose. At their farewell she admits her affection: “Mais oui, je t'aime”---“Yes, of course I love you.” She explains that her behaviour has concealed it, and then urges him to leave rather than watch her cry \citep[chap.~IX]{saintExupery1943}. Her words disclose something he has failed to understand. The immediate exchange offers a clear reason to read an avowal of feeling without treating it as a demand for a new undertaking to stay. This does not settle what his care for her already obliges him to do elsewhere in the story.

\paragraph{P: does the utterance undertake, renew or invoke a commitment?}
A comparatively explicit example comes from Peter Nansen's \emph{Maria}. Near the end of this first-person account of a love affair, Maria asks whether the narrator's changed attitude is durable, rather than merely a return of sexual longing. His reply includes:
\begin{quote}
“Jeg elsker Dig i Lyst og i Nød, Hverdage og Festdage. Du er den, jeg vil dø med, men først vil jeg leve meget, meget længe med Dig.”

“I love you in joy and in adversity, on ordinary days and feast days. You are the one I want to die with, but first I want to live a very, very long time with you.”
\end{quote}
The declaration answers a question about continuing union and uses a vow-like formula extending across circumstances \citep[sec.~LXI]{nansenMaria}. The corpus annotation assigns T4/P4/E0: the avowal participates in an assurance about their life together.\caseref{nansen} Here and below, translations of quoted non-English passages are ours; the original wording and saved analytical translations remain available.

The more familiar \emph{Ghost} dramatizes why merely conveying the same apparent information may feel insufficient. Molly proposes marriage to Sam and asks why he answers her declarations with “ditto”: “You say ‘ditto’. It's not the same.” He replies that people use the words too readily \citep{rubinGhost}. The scene invites an undertaking interpretation of what Molly wants, but does not prove that she is demanding a particular promise. She might also want explicit acknowledgment or emotional openness. This ambiguity is useful: identifying possible P is a question to investigate, not an automatic consequence of a declaration's importance.

\paragraph{E: is producing the expression itself presented as an affective response?}
In \emph{Up}, Carl and Russell meet Dug, a dog whose collar turns his thoughts into speech. Dug introduces himself with “I have just met you and I love you.” He immediately jumps up on Carl; the screenplay describes a device that continuously gives words to his thoughts \citep[pp.~46--47]{docterPetersonUp}. The comic machinery makes affection emerging directly into words particularly easy to imagine. It motivates E without establishing that an ordinary human declaration works the same way. A more direct corpus example below explicitly describes words escaping a speaker's inhibition. Mere excitement, repetition or an exclamation mark would not suffice under our criteria.

The dimensions can combine. An avowal can renew a commitment; an affective outburst can also present a feeling as true. P is narrower than Austin's broad concern with actions accomplished in speaking \citep{austin1975}: an avowal itself is an act, but this does not automatically make it a relational undertaking. E, correspondingly, is narrower than being emotionally expressive. The classifications concern what is presented, what is undertaken, and how the words are represented as being produced.

One distinction is essential throughout: \emph{T is not a sincerity score}. A false avowal can have truth-presenting force. Likewise, a person can undertake something without intending to honour it, and a seemingly spontaneous declaration may be calculated. The presence of an affection, the avowal of that affection, and the obligations arising from the avowal are separate questions. Negation and tense also matter: “I no longer love you” can present a condition as true without asserting present love.

The categories leave substantial work to be done. Under T, what relation is at issue---romantic attachment, familial affection, friendship? Under P, is the act a promise, an offer, an acceptance, or an invocation of an existing obligation? What are its terms? Under E, does the text depict involuntary verbal production, a release of inhibition, or the metaphorical overflow of feeling? We argue for the usefulness of the broad distinctions while retaining these differences within them.

\section{A multilingual corpus with an inspectable procedure}

The initial public corpus comprises eight English works, ten French works, four Norwegian works, and two each in Danish, German, Italian and Swedish. Extraction version 0.11 identified 227 candidates in their complete canonical texts. Automated membership review retained 225, all classified with prompt version 0.3.1. An additional private work in the wider project is excluded from the results reported here. “Public” refers to the availability of the research material: the corpus includes a contemporary licensed novel as well as historical texts \citep{lmcw2026}.

The subsequent extension was selected to examine declarations associated with duty, dependence, marriage and the sexual double standard. It contains Victoria Benedictsson's \emph{Pengar}, Amalie Skram's \emph{Constance Ring} and \emph{Lucie}, George Gissing's \emph{The Odd Women}, and Elizabeth von Arnim's \emph{Vera}. The complete novels were processed, rather than only the motivating scenes. Table~\ref{tab:extension} records the yield, including the zero for \emph{Pengar}. This is a purposive extension, not a representative sample of feminist fiction or a test of differences between genres. The combined corpus contains 35 works, with ten in English, ten in French, six in Norwegian, three in Swedish, and two each in the other three languages.

\begin{table}[htbp]
\centering
\caption{The targeted extension. Support thresholds are descriptive; T=4 and O=0 for all 27 retained occurrences.}
\label{tab:extension}
\begin{tabular}{llrrr}
\toprule
Work (original publication) & Language & Retained & P$\geq2$ & E$\geq2$ \\
\midrule
\emph{Pengar} (1885) & Swedish & 0 & 0 & 0 \\
\emph{Constance Ring} (1885) & Norwegian & 11 & 2 & 0 \\
\emph{The Odd Women} (1893) & English & 12 & 1 & 1 \\
\emph{Vera} (1921) & English & 2 & 1 & 0 \\
\emph{Lucie} (1888) & Norwegian & 2 & 1 & 0 \\
\midrule
Total & & 27 & 5 & 1 \\
\bottomrule
\end{tabular}
\end{table}

The new Norwegian material motivated extraction version 0.12, which adds patterns for intervening \emph{også}/\emph{ogsaa} (“also”) and historical \emph{elskede}. It raises the extension's candidate count from 23 to 27; a comparison records unchanged candidate spans in the original 30 public works. All 27 new candidates were provisionally retained. The two Skram texts derive from National Library of Norway first-edition scans and OCR, with three documented, scan-checked corrections in \emph{Constance Ring}. They remain imperfect OCR texts, not fully proofread editions. The original stage and the extension therefore retain their own extraction and annotation records.

Our units are \emph{ILY-occurrences}. The rules admit negative, past, future and modal forms, embedding, and quoted, imagined or hypothetical speech. Not every retained occurrence is an actual present-tense declaration between characters. Searching complete works gives routine occurrences a place beside memorable ones, but does not guarantee recall of every expression of love. In particular, a paraphrased wedding vow or an unspoken affection may be relevant to the argument without meeting the explicit first-person/second-person extraction rule. Appendix~\ref{app:corpus} gives the work inventory and retained counts.

The saved annotation requests use GPT-5.6 Sol. They supply the matched wording, local passage and wider context, normally approximately 3,000 characters on either side expanded to paragraph boundaries. Non-English inputs include generated analytical translations; the originals remain authoritative. Reliable background knowledge is permitted but must be identified: the recorded outputs mark it as used in two baseline cases and no extension cases. Thus access to complete source texts should not be confused with the annotator having reread a complete novel for each judgment.

The output assigns independent integer support scores from 0 to 4, together with an utterance-status label, analysis, evidence and uncertainty. Scores do not sum to a fixed total and are not calibrated probabilities. The prompt deliberately requires positive evidence for E and places a high burden of proof on O, an additional category for an important core function that no combination of T/P/E adequately represents. Consequently, low O cannot independently establish the framework's completeness, and low E cannot establish that love declarations are seldom emotional.

The advantage of automation is practical revisability. Published criteria, contexts and judgments allow an objection to be turned into a revised protocol and applied across the same material. A human research team can also make its methods explicit; a model's written rationale does not expose every influence on its answer. Here, automation reduces the labour of repeated multilingual application. The saved runs estimate US\$14.50 for the 225 baseline classifications and US\$1.22 for the 27 new classifications. Translation, acquisition, membership review and method development are additional costs. These are recorded estimates, not a complete research budget.

The close readings below are purposive interpretations of the recorded evidence. They include apparent successes, boundary cases and the failure of an expected pattern to appear. They are not independent human validation or an inter-annotator reliability experiment. Historical scores are reported unchanged; where the reading questions a score, the disagreement is explicit. The source register links each discussed scored occurrence to the exact request and output.

\section{What the recorded patterns show}

\begin{table}[htbp]
\centering\small
\caption{Assigned support scores, with the original corpus and targeted extension reported separately.}
\label{tab:scores}
\begin{tabular}{llrrrrr}
\toprule
Stage & Dimension & 0 & 1 & 2 & 3 & 4 \\
\midrule
Original (225) & T & 0 & 4 & 3 & 9 & 209 \\
 & P & 145 & 43 & 16 & 10 & 11 \\
 & E & 206 & 2 & 10 & 6 & 1 \\
 & O & 224 & 0 & 0 & 1 & 0 \\
\midrule
Extension (27) & T & 0 & 0 & 0 & 0 & 27 \\
 & P & 16 & 6 & 3 & 1 & 1 \\
 & E & 26 & 0 & 1 & 0 & 0 \\
 & O & 27 & 0 & 0 & 0 & 0 \\
\midrule
Combined (252) & T & 0 & 4 & 3 & 9 & 236 \\
 & P & 161 & 49 & 19 & 11 & 12 \\
 & E & 232 & 2 & 11 & 6 & 1 \\
 & O & 251 & 0 & 0 & 1 & 0 \\
\bottomrule
\end{tabular}
\end{table}

\begin{table}[htbp]
\centering\small
\caption{Dimensions reaching 2. The direct-speech column is a subset of the combined corpus. O is not used to define these combinations.}
\label{tab:patterns}
\begin{tabular}{lrrrr}
\toprule
Dimensions & Original & Extension & Combined & Direct speech \\
\midrule
T & 168 & 21 & 189 & 149 \\
T and P & 36 & 5 & 41 & 34 \\
T and E & 16 & 1 & 17 & 11 \\
T, P and E & 1 & 0 & 1 & 1 \\
None & 4 & 0 & 4 & 0 \\
\midrule
Total & 225 & 27 & 252 & 195 \\
\bottomrule
\end{tabular}
\end{table}


Table~\ref{tab:scores} keeps the two stages distinguishable. The extension increases the number of retained occurrences from 225 to 252, with five additional P scores of at least 2 and one additional E score of at least 2. Its T scores are uniformly 4. Across both stages, P reaches 2 in 42 records and E in 18. Every one of these records also reaches 2 on T. These judgments describe widespread avowal with variation in undertaking and expressive/reflexive force, rather than a division into three exclusive classes.

Table~\ref{tab:patterns} summarizes combinations at the threshold of 2. The final column restricts the combined material to the 195 records labelled direct speech. Written, imagined, hypothetical and other statuses remain in the all-occurrences columns. Several occurrences can come from one scene, so none of these totals should be read as a count of independent conversations, speakers or relationships.

The threshold is a display convention, not a validated boundary. Raising it to 3 produces two baseline cases that reach the P threshold but not the T threshold: both have T2/P4. One is a conditional marital offer in \emph{The Tenant of Wildfell Hall}; the other offers future love as an inducement in \emph{Madame Bovary}. We discuss them below. It would be misleading to say that the corpus contains no variation in the relative strength of T and P. The narrower negative result is that it contains no record with strong P (3 or 4) and weak T (0 or 1), and that the targeted extension contains no reduced-T case at all.

High T partly reflects the extraction of explicit love constructions and a criterion that includes false avowals. The P and E distributions are likewise conditional on the instructions. Counts show where interpretations have been assigned systematically; whether the distinctions are useful depends on the textual warrants. The following passages make those warrants, and their limitations, available for inspection.

\section{Avowal and undertaking come apart in the texts}

\subsection{Affection while withholding the requested commitment}

Charlotte Bront\"e's Jane has been asked by St John Rivers to marry him and accompany his missionary work in India. She is willing to help him in a sisterly role, but rejects marriage. When he accuses her of retreating from an engagement, she answers:
\begin{quote}
“There is no dishonour, no breach of promise, no desertion in the case. I am not under the slightest obligation to go to India, especially with strangers. With you I would have ventured much, because I admire, confide in, and, as a sister, I love you.”
\end{quote}
The passage directly connects an avowal with a dispute about obligation \citep[chap.~XXXV]{bronteJaneEyre}. The T4/P0/E0 annotation is supported by Jane's assertion of affection and explicit denial of the undertaking attributed to her.\caseref{jane} “As a sister” identifies the relationship being avowed: the broad T label alone would not distinguish this affection from romantic attachment. Jane's words do not deny all obligations to St John; they reject the particular inference he is pressing.

The new Gissing material supplies an especially useful romantic comparison. Rhoda Nunn, who works for women's independence, is considering Everard Barfoot's proposal of a union without formal marriage. He presses for an immediate decision. She asks for time, and he asks whether she doubts her love:
\begin{quote}
“No. If I understand what love means, I love you.”
\end{quote}
She nevertheless withholds the requested kiss and removes the marriage ring he puts on her finger. Later in the same scene she accepts the prospect of a civil marriage \citep[chap.~XXV]{gissingOddWomen}. At the point classified T4/P0/E0, however, the avowal has not supplied the assent Everard demands.\caseref{rhoda} The conditional phrase qualifies her understanding of love, not her readiness to say that she loves him. The narrator explicitly represents passionate feeling. Love's presence and the terms on which she will act upon it remain distinguishable even though the negotiation subsequently advances.

George Eliot makes the disputed inference almost an explicit argument. Fred Vincy, whose idleness and debts trouble Mary Garth, wants her to declare love and promise marriage. Mary replies: “If I did love you, I would not marry you: I would certainly not promise ever to marry you” \citep[chap.~XIV]{eliotMiddlemarch}. This hypothetical occurrence receives T1/P0/E0.\caseref{mary} It does not assert Mary's love, but it directly disputes Fred's assumption that love would require a marital undertaking. Its value to the argument is not interchangeable with that of an enacted avowal.

Ibsen brings us back to the opening problem of changed feelings. After the crisis over Nora's secret loan, Torvald first condemns her and then expects the marriage to resume when the danger to himself passes. In their final conversation Nora says: “Men jeg kan ikke gjøre ved det. Jeg elsker deg ikke mer”---“But I cannot help it. I do not love you any more” \citep[act~III]{ibsenDukkehjem}. The T4/P0/E0 annotation treats the words as presenting a changed condition, while the surrounding exchange articulates departure.\caseref{nora} The reading does not resolve the play's ethical argument. It identifies why an accusation of breaking a promise cannot simply be read off the love clause. The quotation follows the modernized Norwegian text used in the corpus.

\subsection{Offering oneself and accepting a relationship}

In \emph{Constance Ring}, Lorck declares himself to Constance while she is still married to Ring. The scene is unwelcome: she rejects his advances, asks him not to frighten her, and tries to withdraw. He persists:
\begin{quote}
“Jeg elsker dig, Constance \ldots\ Jeg giver dig mig selv, mit Liv, min Sjæl,---vil du ikke ta imod det?”

“I love you, Constance \ldots\ I give you myself, my life, my soul---will you not accept it?”
\end{quote}
The offer and request for acceptance explain the T4/P4/E0 classification \citep[p.~137]{skramConstance}.\caseref{lorck} They also expose a limitation of the P label. An attempted undertaking need not be welcome, accepted or ethically admirable. Lorck's threats and physical pressure make that distinction particularly important. His offer does not create a reciprocal commitment on Constance's behalf. The strong P score is also open to a narrower reading: how much undertaking force belongs to the love clause itself, and how much to the explicit offer that follows?

In \emph{Lucie}, Skram supplies a contrasting exchange in which a woman accepts. Henny tests Knut's sexual double standard by claiming to have had lovers. After jealousy and distress, he says that he would still beg her to take him as her husband. She answers:
\begin{quote}
“Og jeg elsker også Dig. Jeg tror på Dig og jeg er din. Siden Du vil ha mig.”

“And I love you too. I believe in you and I am yours. Since you want me.”
\end{quote}
She subsequently describes herself as an engaged woman and reveals that the supposed lovers were invented \citep[pp.~176--177]{skramLucie}. Her T4/P3/E0 occurrence participates in accepting the proposed relationship.\caseref{henny} The explicit “I am yours” and her own account of the engagement give P more definite support than mere affectionate delivery would. The passage concerns the conditions of acceptance; it is not evidence of a loveless undertaking.

Two more moderate judgments help locate the interpretive boundary. In \emph{The Odd Women}, Widdowson courts Monica, a young woman with poor employment prospects. She insists that she remain free. His answer joins “I love you with all my soul” to “your rules shall be mine” and a request for a chance to win her; the record is T4/P2/E0 \citep[chap.~VII]{gissingOddWomen}.\caseref{widdowson} Lorck's later, carefully composed letter to Constance also receives T4/P2/E0. After prolonged reflection, he writes, he knows “til Bunden af min Sjæl, at jeg elsker Dem”---“to the depths of my soul that I love you” \citep[p.~309]{skramConstance}.\caseref{lorckletter} Both are serious overtures, but their undertaking force is less explicit than Henny's acceptance. The letter is especially debatable under a prompt that says courtship alone does not establish P. These scores identify questions for close reading; they do not turn courtship automatically into promise.

\subsection{Commitment without an avowal of present love?}

Two baseline cases already complicate the dominance of T. In Anne Bront\"e's \emph{The Tenant of Wildfell Hall}, Helen recounts refusing Mr Boarham's proposal. He interrupts her refusal: “I would love you, cherish you, protect you” \citep[chap.~XVI]{bronteTenant}. The T2/P4/E0 classification treats this mainly as a conditional marital offer, with some presentation of prospective feeling.\caseref{boarham} In \emph{Madame Bovary}, Emma, desperate for money, urges L\'eon to obtain it: “Oh! tâche! tâche! je t'aimerai bien!”---“Oh, try! Try! I shall love you dearly!” \citep[part~III, chap.~VII]{flaubertBovary}. The same scores describe a promise-like inducement; the expression may also suggest the conventional “I'll love you for it.”\caseref{emmafuture} Neither is a straightforward report of present affection. Both nevertheless retain some T support under a criterion that includes prospective states.

Tolstoy supplies a different, unscored comparison. In \emph{War and Peace}, Pierre scarcely knows how to resist the expectations arranging his marriage to H\'el\`ene. Alone with her, he remembers the appropriate social formula and says “Je vous aime.” Six weeks later they are married. After the duel with Dolokhov, he recalls the declaration as “a lie, and worse than a lie” \citep[book~III, chap.~II; book~IV, chap.~VI]{tolstoyWarPeace}, directly linking the words to binding himself to her.

This motivates investigating undertaking force when affection is absent or uncertain, but does not establish an absence of truth-presenting force. Calling the declaration a lie is compatible with acknowledging that force; “worse” can indicate binding consequences, self-deception or damage. These are interpretations of Pierre's retrospective judgment, not Tolstoy's explicit semantic gloss. The earlier scene also has Pierre telling himself that he loves her; the difference between his accounts matters. High P with absent \emph{felt affection} is distinct from high P with absent \emph{T as defined here}.

\section{What the targeted extension teaches us}

\subsection{A solicited avowal and the work of self-persuasion}

The \emph{Constance Ring} passage that initially attracted attention is not Lorck's high-P declaration just discussed. It occurs much later, during Constance's second marriage, to Lorck. They recall their happiness together and anticipate their child. Lorck asks her to say that she loves him. She briefly wants to shake off his embrace and flee, remembering the pleasure she has felt in another man's proximity. She then deliberately returns to the memories of happiness that have sustained the marriage. The narration describes her wrapping those memories around herself as protection. She says:
\begin{quote}
“Jeg elsker dig \ldots\ elsker, elsker dig”---“I love you \ldots\ love, love you.”
\end{quote}
The kisses and caresses that follow are exceptionally ardent \citep[chap.~XXIV, pp.~369--370]{skramConstance}. The annotation is T4/P0/E0.\caseref{conscience} It treats the words as a solicited assertion reinforced by self-persuasion, without clear evidence of a fresh or renewed undertaking. The deliberate effort preceding them also weighs against reading their intensity alone as E.

The narration prevents an uncomplicated reading of transparent feeling. It does not establish that Constance has never loved Lorck: her remembered happiness is part of the evidence. Nor does saying something to sustain a marriage necessarily make the utterance a promise. One might argue that compliance with Lorck's request reaffirms their bond and warrants a higher P score. To make that argument persuasive, we must specify what is being undertaken beyond the reassurance. The case sharpens the question with which the extension began, rather than straightforwardly confirming our initial expectation.

\subsection{Private vows alongside spoken love}

In \emph{Vera}, Lucy is newly married to Wemyss, whose first wife died after a fall at his house. Standing with him at the window, she is distressed by the terrace where Vera died. He treats the view as something to admire. Lucy tries to believe that his untroubled manner is the right response, seeks his kisses, and says:
\begin{quote}
“Oh, I love you, love you----”
\end{quote}
At the same time, she makes “secret vows of sensibleness, of wholesomeness, of a determined, unfailing future simplicity” \citep[chap.~XVII]{vonArnimVera}. The surrounding narration follows Wemyss's proprietary thoughts about her face and his unrestricted right to kiss it. The record assigns T4/P2/E0.\caseref{lucy}

Affection and self-discipline are intertwined here. The explicit vows support the interpretation of commitment \emph{in the scene}, but their secrecy creates a difficulty for locating it in the spoken ILY. Wemyss cannot accept terms that she has not communicated. The saved annotation itself allows that P could be lower if the mental act and spoken avowal are separated. This is a useful challenge to consistency: in some other records an adjacent marriage proposal is distinguished from the love clause, whereas here simultaneous private resolutions contribute to P. The appropriate unit of interpretation---clause, speech turn, or whole scene---needs clearer treatment before these intermediate scores can bear much numerical weight.

\subsection{A relevant scene with no extracted occurrence}

\emph{Pengar} returns no eligible match. It nevertheless contains one of the extension's most pertinent passages. Selma, married while very young to a financially secure older man, reads his habitual affectionate formulas in a letter. She wonders whether she has ever said she loves him. She remembers the wedding formula, which she repeated “själlöst som en papegoja”---“soullessly, like a parrot” \citep[chap.~VI]{benedictssonPengar}.

The scene separates institutional commitment from a remembered experience of love. However, it recounts and reflects on a formula rather than supplying an eligible explicit ILY token. Its evidential role is therefore qualitative; it must not be silently added to the 27 classifications. A similar limitation applies to the motivating \emph{Odd Women} scene in which Widdowson demands “Say you love me!” The narrator reports: “She kissed his cheek, but did not utter a word” \citep[chap.~XVI]{gissingOddWomen}. Refusal or absence can be significant precisely because the expected words do not occur. The extraction boundary is a methodological choice with interpretive consequences.

The extension thus supplies additional contrasts and exposes limitations of both search and classification. It yields two strong P cases (3 or 4), three moderate P cases (2), and no low-T/high-P case. This does not show that declarations made from duty are nonexistent or insignificant. It shows that selecting works with those situations did not produce the anticipated score pattern under this protocol. A future account should track the represented presence of affection separately from the force of avowing it, and distinguish an interpersonal undertaking from an effect such as reassurance or self-persuasion. These should be explicit additions to the analysis, not retrospective changes to the meaning of T.

\section{Expression, narrative perspective and the limits of scoring}

\subsection{When the text describes the words escaping}

In Ania Cofield's contemporary romance \emph{Nikki's Touch}, Nikki and Keith face an approaching separation after an intimate encounter. Lying together afterwards, she says “I love you,” with the narration specifying that the words are “spilling out before she could stop them” \citep{cofieldNikkisTouch}. The T4/P0/E4 classification has direct evidence about verbal production preceding inhibition.\caseref{nikki} The narrative also presents the admission as honest and releasing. E accompanies T; it does not cancel the presentation of a feeling as true. Nothing in that combination alone determines what Nikki subsequently owes Keith.

The contrast with \emph{Persuasion} is instructive. Anne Elliot and Frederick Wentworth were engaged more than eight years earlier before she was persuaded to break the engagement. Having returned, Wentworth overhears Anne defending women's constancy and writes a renewed offer. His letter combines “I have loved none but you” with “I offer myself to you again” \citep[chap.~XXIII]{austenPersuasion}. The love occurrence receives T4/P2/E0.\caseref{wentworth} His difficulty in writing conveys great emotion, but gives less direct evidence than Nikki's passage that the target words escaped control. It also illustrates again the question of how the love clause contributes to an undertaking explicitly made elsewhere in the same turn.

Flaubert's Emma gives a less determinate E case. During her affair with Rodolphe, her recurrent declarations include “je t'aime à ne pouvoir me passer de toi”---“I love you so much that I cannot do without you.” Rodolphe discounts the conventional language, while the narrator warns against assuming that familiar words cannot carry deeply felt emotion \citep[part~II, chap.~XII]{flaubertBovary}. Two occurrences receive T4/P0/E3, with the annotations invoking the narrator's image of feeling overflowing.\caseref{emmaexpressive} The metaphor supports expression more directly than it establishes involuntary verbal production. The difference from Nikki's passage is worth retaining within E.

Nor should bodily expression be conflated with speaking. In another \emph{Jane Eyre} passage, St John's religious vocation prevents him from pursuing Rosamond. Jane reads his involuntary response to her as a declaration: “He seemed to say, with his sad and resolute look, if he did not say it with his lips, ‘I love you, and I know you prefer me.’” The record is labelled nonverbal verbalized, with T4/P0/E3 \citep[chap.~XXXII]{bronteJaneEyre}.\caseref{stjohn} It concerns narration translating a look into words. That can illuminate affective expression, but it is not an instance of words reflexively uttered by St John. The aggregate E totals include such differences in the events being classified.

\subsection{Neither sincerity nor commitment is settled by a score}

The opening love test in \emph{King Lear} provides a familiar unscored reminder. Lear makes his daughters' shares depend on declarations of love. Cordelia refuses the competition in flattery, but she does not simply deny her affection: “I love your Majesty / According to my bond, no more nor less” \citep[1.1]{shakespeareLear}. Her subsequent loyalty makes the sincerity of that affection central to the tragedy. It is compatible with resisting a demanded verbal performance. The sisters' flattering claims, conversely, would not cease to have truth-presenting form merely because their sincerity is doubtful. This distinction helps prevent T from becoming a measure of whether a character truly loves.

Lawrence explicitly makes the formula's adequacy problematic. In \emph{Women in Love}, Birkin and Ursula are approaching marriage, but their interpretations of their intimacy diverge: “Even when he said, whispering with truth, ‘I love you, I love you,’ it was not the real truth.” The narrator explains that the conventional expression fails to capture Birkin's transformed sense of self \citep[chap.~XXVII]{lawrenceWomenInLove}. The target receives T3/P1/E0/O3; a nearby declaration receives T4/P3/E0/O0.\caseref{birkin} The targets differ, so their scores are not automatically contradictory. They nevertheless demand a principled explanation of why a contextual qualification counts as pressure against the ontology in one case and not the other. The solitary positive O classification is an invitation to examine that distinction, not evidence that all other complexity has been resolved.

\section{Discussion: what has been established, and what remains open}

The strongest result is a set of textual contrasts that a uniform account would need to explain. Jane avows sisterly love while explicitly denying an alleged engagement. Rhoda avows romantic love while withholding the requested assent. Nansen's declaration participates in a vow-like assurance; Henny's in an acceptance. Nikki's narration identifies a failure of inhibition that an account only of propositional content or deliberate undertaking would omit. The corpus makes these cases part of a systematic collection, and the recorded rationales show how the same criteria have been applied to less memorable occurrences.

The evidence is strongest for variation in P alongside avowal, with more heterogeneous evidence for E. It does not establish three independent, exclusive classes. Nor does it establish their exhaustiveness. In particular, the absence of a low-T/high-P outcome in the targeted extension must be reported as a limit of the result. The broader definition of T, which includes false or self-deceptive avowals, makes that finding compatible with the very situations of duty and dependence that motivated the search.

The ethical importance survives these qualifications. When people disagree over an earlier declaration, the issue may be whether an undertaking occurred, what it concerned, or whether the claimed feeling was present. A description of words escaping inhibition raises another question, without automatically excusing their consequences. Intention, uptake, coercion, reliance and subsequent conduct can pull in different directions. Lorck's forceful offering of himself does not oblige Constance to accept it. Lucy's private resolution does not automatically become an undertaking to Wemyss. A single summary score cannot adjudicate all these perspectives.

An idealized romance may align sincere affection, freely made commitment and an apparently spontaneous declaration. \citet{rayner2014} proposes \emph{Splash} as an illustration of that attractive convergence. It suggests a genre hypothesis: do some romances more often present their central declarations as moments when feeling, commitment and expression agree? Our current material cannot answer it. One contemporary romance is not a genre sample, and emotional intensity or apparent ease of expression need not satisfy the restrictive E criterion. A convincing test would require a separately defined comparison corpus and an explicit account of what counts as convergence, including where in a story a declaration occurs.

Further work should first address the concrete uncertainties the passages expose. How much of the surrounding turn belongs to the target's force? How should asserted, hypothetical and nonverbal occurrences be compared? Does E require evidence about the production of words, or can a narrator's metaphor of emotional overflow suffice? Independent readings and controlled variations in context and instructions could test these questions. The selection and pattern development are exploratory, and neither one model run nor these subsequent readings supply a gold standard.

There is nevertheless a useful methodological gain. Objections can be attached to specific wording, contexts, criteria and records, and revised criteria can be reapplied without reconstructing the corpus. The contribution is this revisable body of evidence together with a first account of its interpretive contrasts. It gives firmer grounds for asking what has been avowed, undertaken or expressed when familiar words become the basis of conflicting expectations.

\paragraph{Data and drafting history.}
The analyses use repository snapshot \href{https://github.com/mannyrayner/let_me_count_the_ways/tree/ef7ea7933d8aa02b176d8fbddc496e55a41a7144}{\texttt{ef7ea79}}. The two annotation runs, full prompts, supplied contexts, source provenance and outputs remain separately identifiable. Draft sources, a record of revisions, the case register and the script deriving the tables accompany this manuscript under \texttt{docs/paper/humanities/}. Introductory illustrations and the Tolstoy and Lear comparisons have not been added to the scored dataset. The publicly available membership decisions remain provisional AI judgments pending human review.

\bibliographystyle{plainnat}
\bibliography{references}

\clearpage
\appendix
\section{Corpus inventory}
\label{app:corpus}
B denotes the original stage; X denotes the targeted extension. Language codes: da Danish, de German, en English, fr French, it Italian, no Norwegian, sv Swedish. Historical Dano-Norwegian is coded no, consistently with the original corpus. Retained counts reproduce provisional AI membership decisions.

\begingroup\small
\setlength{\tabcolsep}{3pt}
\begin{longtable}{@{}p{34mm}p{71mm}ccr@{}}
\caption{The 35 complete works searched, including the zero-yield extension work.}\\
\toprule
Author & Work & Lang. & Stage & Retained \\
\midrule\endfirsthead
\toprule
Author & Work & Lang. & Stage & Retained \\
\midrule\endhead
\midrule\multicolumn{5}{r}{Continued on next page}\\\endfoot
\bottomrule\endlastfoot
Herman Bang & \emph{Ved Vejen} & da & B & 1 \\
Peter Nansen & \emph{Maria: En Bog om Kærlighed} & da & B & 2 \\
Theodor Fontane & \emph{Effi Briest} & de & B & 2 \\
Johann Wolfgang von Goethe & \emph{Die Leiden des jungen Werther} & de & B & 1 \\
Louisa May Alcott & \emph{Little Women} & en & B & 5 \\
Jane Austen & \emph{Persuasion} & en & B & 1 \\
Charlotte Brontë & \emph{Jane Eyre} & en & B & 9 \\
Anne Brontë & \emph{The Tenant of Wildfell Hall} & en & B & 10 \\
Ania Cofield & \emph{Nikki's Touch} & en & B & 11 \\
George Eliot & \emph{Middlemarch} & en & B & 5 \\
D. H. Lawrence & \emph{Women in Love} & en & B & 11 \\
Edith Wharton & \emph{The Age of Innocence} & en & B & 2 \\
Honoré de Balzac & \emph{Illusions perdues} & fr & B & 11 \\
Colette & \emph{Le blé en herbe} & fr & B & 4 \\
Benjamin Constant & \emph{Adolphe} & fr & B & 1 \\
Alexandre Dumas fils & \emph{La Dame aux camélias} & fr & B & 19 \\
Gustave Flaubert & \emph{Madame Bovary} & fr & B & 15 \\
Madame de La Fayette & \emph{La Princesse de Clèves} & fr & B & 8 \\
Edmond Rostand & \emph{Cyrano de Bergerac} & fr & B & 20 \\
George Sand & \emph{La Mare au Diable} & fr & B & 5 \\
Madame de Staël & \emph{Corinne; ou, l'Italie} & fr & B & 9 \\
Stendhal & \emph{Le Rouge et le Noir} & fr & B & 11 \\
Guido da Verona & \emph{Colei che non si deve amare} & it & B & 16 \\
Luciano Zùccoli & \emph{L'amore di Loredana} & it & B & 12 \\
Knut Hamsun & \emph{Pan} & no & B & 4 \\
Knut Hamsun & \emph{Victoria} & no & B & 11 \\
Henrik Ibsen & \emph{Et dukkehjem} & no & B & 3 \\
Sigrid Undset & \emph{Kristin Lavransdatter} & no & B & 7 \\
Selma Lagerlöf & \emph{Gösta Berlings saga} & sv & B & 7 \\
August Strindberg & \emph{Fröken Julie} & sv & B & 2 \\
George Gissing & \emph{The Odd Women} & en & X & 12 \\
Elizabeth von Arnim & \emph{Vera} & en & X & 2 \\
Amalie Skram & \emph{Constance Ring} & no & X & 11 \\
Amalie Skram & \emph{Lucie} & no & X & 2 \\
Victoria Benedictsson & \emph{Pengar} & sv & X & 0 \\
\midrule
\multicolumn{4}{r}{Total} & 252\\
\end{longtable}
\endgroup


\ifincludesupplement
\clearpage
% Optional supplement, included by let_me_count_the_ways.tex.
\section{Editorial and source notes}

This supplement preserves the audit trail for working draft v0.3. It is not a new scored dataset or a completed context-dependence audit. The main paper and these notes use repository snapshot \nolinkurl{662f77a398e884f9bf30601c3bbf433a722c9619}. The private thirty-first work is outside the processed snapshot. Historical annotations have not been changed.

\subsection{Revision and remaining questions}

The opening and discussion now centre on the obligations attributed to a love declaration. T/P/E is presented as a first level of analysis, with further questions about the loving relation, precise illocutionary force and mode of expressive production. The examples have been reorganized around those distinctions. The paper does not infer moral innocence or wrongdoing from a score alone.

The joint-distribution table is newly calculated from the unchanged frozen summary. Its direct-speech column is a restriction to an existing event-status label, not a reclassification. The reading register below is unchanged. The theoretical passage remains attributed to Barthes with its scope stated; the author of \emph{Nikki's Touch} remains corrected to Ania Cofield.

The following questions remain for a later revision:
\begin{enumerate}
\item Develop the scholarly positioning beyond the initial Austin and Barthes discussion, including relevant work on avowal, commitment and literary speech. Avoid treating truth-conditional semantics as if it necessarily excluded pragmatics. The current discussion is not yet a full literature review. The proposed learning-from-examples account is a motivating possibility, not a result of this corpus study.
\item Clarify how scores for asserted, embedded, imagined and nonverbal cases should be compared. The existing event-status field combines medium, embedding and event realization. A supplementary description may suffice; reannotation should not be assumed necessary in advance.
\item Distinguish explicit verbal-production evidence from bodily expression and metaphorical overflow. Preserve disagreement where the text does not decide it, and retain historical scores if revised judgments are later added.
\item Obtain human feedback on the central readings and source-language interpretation. Current close reading is purposive and does not constitute independent human validation. The bibliography identifies the electronic texts used; a later submission may also require critical-edition references.
\end{enumerate}

\subsection{Repository records}

The following links point to the frozen snapshot, rather than the current repository head.
\begin{itemize}
\item \href{https://github.com/mannyrayner/let_me_count_the_ways/tree/662f77a398e884f9bf30601c3bbf433a722c9619/corpus/works/}{Corpus work manifests and canonical texts}.
\item \href{https://github.com/mannyrayner/let_me_count_the_ways/blob/662f77a398e884f9bf30601c3bbf433a722c9619/results/review/canonical_31_v0_11_ai_review_v1/summary.json}{Membership-review summary}.
\item \href{https://github.com/mannyrayner/let_me_count_the_ways/blob/662f77a398e884f9bf30601c3bbf433a722c9619/results/annotation/canonical_31_v0_11_v0_3_1/summary.json}{Annotation summary}.
\item \href{https://github.com/mannyrayner/let_me_count_the_ways/blob/662f77a398e884f9bf30601c3bbf433a722c9619/prompts/annotation/classify_passage_v0_3_1.md}{Annotation prompt, version 0.3.1}.
\item \href{https://github.com/mannyrayner/let_me_count_the_ways/blob/662f77a398e884f9bf30601c3bbf433a722c9619/prompts/annotation/classification_schema_v0_3.json}{Classification schema, version 0.3}.
\item \href{https://github.com/mannyrayner/let_me_count_the_ways/tree/662f77a398e884f9bf30601c3bbf433a722c9619/results/annotation/canonical_31_v0_11_v0_3_1/annotations/}{All saved annotation records}.
\end{itemize}

The request files linked below contain the target expression, local text and wider context actually supplied to the annotator. The output files contain the corresponding score, rationale, evidence and uncertainty. Each occurrence identifier maps unambiguously to these records.

\subsection{Derivation of the joint counts}

Table~\ref{tab:patterns} is calculated from the \texttt{cases} array in the linked annotation summary. For each record, retain exactly those dimensions whose integer score is at least 2. Count the resulting combinations over all 225 records and, separately, over records with \texttt{utterance\_status} equal to \texttt{direct}. The five nonzero patterns are T (168; 130 direct), TP (36; 30 direct), TE (16; 11 direct), TPE (1; 1 direct), and no dimension reaching threshold (4; 0 direct). These sum to 225 and 172 respectively. Other combinations have zero counts at this threshold.

Distinct-work and distinct-language totals are calculated using the existing \texttt{work\_id} and \texttt{language} fields among records meeting the same threshold. P has 37 records in 17 works and four languages; E has 17 records in ten works and five languages. This is a summary of assigned scores, with no new annotation calls or altered judgments. Threshold sensitivity has not been evaluated in this revision.

The cost figure in the methods section rounds the frozen summary's \texttt{estimated\_total\_cost\_usd} value, 14.502784, to US\$14.50. It describes the run's estimate for the 225 classification calls, not an independently verified invoice or the full cost of the project.

\section{Exploratory reading register}

The first group followed the handover and selected contrasts in score or utterance status. Ordinary direct T4/P0/E0 examples were then chosen by stable hash ordering within languages where such examples existed. Ibsen and further E cases broadened the comparisons; two Nansen records and a further Anne Bront\"e record were added during reading. The 28 records span 15 works and all seven languages. The two Nansen records and two of the Flaubert records share source passages, so the records do not represent 28 independent scenes.

Scores are ordered T/P/E/O. Language codes are da (Danish), de (German), en (English), fr (French), it (Italian), no (Norwegian) and sv (Swedish). Event statuses reproduce the recorded labels, with spaces replacing underscores for readability. ``Input'' and ``output'' link to the saved request and annotation respectively.

\begingroup
\footnotesize
\setlength{\tabcolsep}{4pt}
\renewcommand{\arraystretch}{1.15}
\begin{longtable}{@{}>{\raggedright\arraybackslash}p{66mm}c>{\raggedright\arraybackslash}p{16mm}>{\raggedright\arraybackslash}p{29mm}>{\raggedright\arraybackslash}p{16mm}@{}}
\caption{Selected annotation records retained from draft v0.1.}\\
\toprule
Occurrence & Lang. & T/P/E/O & Event status & Records \\
\midrule
\endfirsthead
\multicolumn{5}{@{}l}{\emph{Table continued}}\\
\toprule
Occurrence & Lang. & T/P/E/O & Event status & Records \\
\midrule
\endhead
\midrule
\multicolumn{5}{r@{}}{\emph{Continued on the next page}}\\
\endfoot
\bottomrule
\endlastfoot
\nolinkurl{austen-persuasion-473794649dc5ac56} & en & 4/2/0/0 & written & \href{https://github.com/mannyrayner/let_me_count_the_ways/blob/662f77a398e884f9bf30601c3bbf433a722c9619/results/annotation/canonical_31_v0_11_v0_3_1/annotations/austen-persuasion-473794649dc5ac56/da01e0aca4075dff2d35f51a87ae5ee7389719960022a5a19e050890f8763c33/request.json}{input}; \href{https://github.com/mannyrayner/let_me_count_the_ways/blob/662f77a398e884f9bf30601c3bbf433a722c9619/results/annotation/canonical_31_v0_11_v0_3_1/annotations/austen-persuasion-473794649dc5ac56/da01e0aca4075dff2d35f51a87ae5ee7389719960022a5a19e050890f8763c33/output.json}{output} \\
\nolinkurl{bronte-jane-eyre-8d223e8c973ddc84} & en & 4/0/3/0 & nonverbal verbalised & \href{https://github.com/mannyrayner/let_me_count_the_ways/blob/662f77a398e884f9bf30601c3bbf433a722c9619/results/annotation/canonical_31_v0_11_v0_3_1/annotations/bronte-jane-eyre-8d223e8c973ddc84/afa503382c40d1751604ad37af8fa33bf0cfd99cc42fbfe2674025d437bd27b2/request.json}{input}; \href{https://github.com/mannyrayner/let_me_count_the_ways/blob/662f77a398e884f9bf30601c3bbf433a722c9619/results/annotation/canonical_31_v0_11_v0_3_1/annotations/bronte-jane-eyre-8d223e8c973ddc84/afa503382c40d1751604ad37af8fa33bf0cfd99cc42fbfe2674025d437bd27b2/output.json}{output} \\
\nolinkurl{cofield-nikkis-touch-be393a6a3d4f0015} & en & 4/0/4/0 & direct & \href{https://github.com/mannyrayner/let_me_count_the_ways/blob/662f77a398e884f9bf30601c3bbf433a722c9619/results/annotation/canonical_31_v0_11_v0_3_1/annotations/cofield-nikkis-touch-be393a6a3d4f0015/8748d81cfb9633467c173fee77c0531c152cc38672ab0895b0905c7fcc4efccd/request.json}{input}; \href{https://github.com/mannyrayner/let_me_count_the_ways/blob/662f77a398e884f9bf30601c3bbf433a722c9619/results/annotation/canonical_31_v0_11_v0_3_1/annotations/cofield-nikkis-touch-be393a6a3d4f0015/8748d81cfb9633467c173fee77c0531c152cc38672ab0895b0905c7fcc4efccd/output.json}{output} \\
\nolinkurl{flaubert-madame-bovary-9c9d14e73da47022} & fr & 4/0/3/0 & quoted or revoiced & \href{https://github.com/mannyrayner/let_me_count_the_ways/blob/662f77a398e884f9bf30601c3bbf433a722c9619/results/annotation/canonical_31_v0_11_v0_3_1/annotations/flaubert-madame-bovary-9c9d14e73da47022/77d181c48ecebf02bc5d8daa7fa51177acb5ea992f5c08cf53a04eb6fb55d7e1/request.json}{input}; \href{https://github.com/mannyrayner/let_me_count_the_ways/blob/662f77a398e884f9bf30601c3bbf433a722c9619/results/annotation/canonical_31_v0_11_v0_3_1/annotations/flaubert-madame-bovary-9c9d14e73da47022/77d181c48ecebf02bc5d8daa7fa51177acb5ea992f5c08cf53a04eb6fb55d7e1/output.json}{output} \\
\nolinkurl{lawrence-women-in-love-86dea678c8044a4f} & en & 4/3/0/0 & direct & \href{https://github.com/mannyrayner/let_me_count_the_ways/blob/662f77a398e884f9bf30601c3bbf433a722c9619/results/annotation/canonical_31_v0_11_v0_3_1/annotations/lawrence-women-in-love-86dea678c8044a4f/c4dfafbe453c344111cfa4a80a4c089ae7bea503bb8690dccd520e6bd201bcbd/request.json}{input}; \href{https://github.com/mannyrayner/let_me_count_the_ways/blob/662f77a398e884f9bf30601c3bbf433a722c9619/results/annotation/canonical_31_v0_11_v0_3_1/annotations/lawrence-women-in-love-86dea678c8044a4f/c4dfafbe453c344111cfa4a80a4c089ae7bea503bb8690dccd520e6bd201bcbd/output.json}{output} \\
\nolinkurl{lawrence-women-in-love-96fc35925c56519b} & en & 3/1/0/3 & direct & \href{https://github.com/mannyrayner/let_me_count_the_ways/blob/662f77a398e884f9bf30601c3bbf433a722c9619/results/annotation/canonical_31_v0_11_v0_3_1/annotations/lawrence-women-in-love-96fc35925c56519b/2bd6ba5b9158b1b4ede8ee93276693e77a37b8ed38f16482cfb232d807c38bab/request.json}{input}; \href{https://github.com/mannyrayner/let_me_count_the_ways/blob/662f77a398e884f9bf30601c3bbf433a722c9619/results/annotation/canonical_31_v0_11_v0_3_1/annotations/lawrence-women-in-love-96fc35925c56519b/2bd6ba5b9158b1b4ede8ee93276693e77a37b8ed38f16482cfb232d807c38bab/output.json}{output} \\
\nolinkurl{eliot-middlemarch-da9a5c10fa10ae23} & en & 1/0/0/0 & hypothetical & \href{https://github.com/mannyrayner/let_me_count_the_ways/blob/662f77a398e884f9bf30601c3bbf433a722c9619/results/annotation/canonical_31_v0_11_v0_3_1/annotations/eliot-middlemarch-da9a5c10fa10ae23/90b4631af4c322e9416dd4dbc176c9b9b841c9d56498d9e7a4fabcc3f8120775/request.json}{input}; \href{https://github.com/mannyrayner/let_me_count_the_ways/blob/662f77a398e884f9bf30601c3bbf433a722c9619/results/annotation/canonical_31_v0_11_v0_3_1/annotations/eliot-middlemarch-da9a5c10fa10ae23/90b4631af4c322e9416dd4dbc176c9b9b841c9d56498d9e7a4fabcc3f8120775/output.json}{output} \\
\nolinkurl{bronte-tenant-of-wildfell-hall-0e20a0f36fe0556a} & en & 2/4/0/0 & quoted or revoiced & \href{https://github.com/mannyrayner/let_me_count_the_ways/blob/662f77a398e884f9bf30601c3bbf433a722c9619/results/annotation/canonical_31_v0_11_v0_3_1/annotations/bronte-tenant-of-wildfell-hall-0e20a0f36fe0556a/f4aba1105ded751735e4ce983961902bbfe5e3e0f182daf0a6a483114d72d5fc/request.json}{input}; \href{https://github.com/mannyrayner/let_me_count_the_ways/blob/662f77a398e884f9bf30601c3bbf433a722c9619/results/annotation/canonical_31_v0_11_v0_3_1/annotations/bronte-tenant-of-wildfell-hall-0e20a0f36fe0556a/f4aba1105ded751735e4ce983961902bbfe5e3e0f182daf0a6a483114d72d5fc/output.json}{output} \\
\nolinkurl{flaubert-madame-bovary-a6b28fb94009e312} & fr & 3/4/0/0 & direct & \href{https://github.com/mannyrayner/let_me_count_the_ways/blob/662f77a398e884f9bf30601c3bbf433a722c9619/results/annotation/canonical_31_v0_11_v0_3_1/annotations/flaubert-madame-bovary-a6b28fb94009e312/c75c1ea81ebaaae946502e5ad0035c53e87f9d0149a730849270addd475301de/request.json}{input}; \href{https://github.com/mannyrayner/let_me_count_the_ways/blob/662f77a398e884f9bf30601c3bbf433a722c9619/results/annotation/canonical_31_v0_11_v0_3_1/annotations/flaubert-madame-bovary-a6b28fb94009e312/c75c1ea81ebaaae946502e5ad0035c53e87f9d0149a730849270addd475301de/output.json}{output} \\
\nolinkurl{bronte-jane-eyre-47ba1ddeb48659a8} & en & 3/4/0/0 & imagined & \href{https://github.com/mannyrayner/let_me_count_the_ways/blob/662f77a398e884f9bf30601c3bbf433a722c9619/results/annotation/canonical_31_v0_11_v0_3_1/annotations/bronte-jane-eyre-47ba1ddeb48659a8/b103c56b9d31f0d69a150b86f9aeb54878086ca3b1da0edcbf1f57eb4bb6e6fb/request.json}{input}; \href{https://github.com/mannyrayner/let_me_count_the_ways/blob/662f77a398e884f9bf30601c3bbf433a722c9619/results/annotation/canonical_31_v0_11_v0_3_1/annotations/bronte-jane-eyre-47ba1ddeb48659a8/b103c56b9d31f0d69a150b86f9aeb54878086ca3b1da0edcbf1f57eb4bb6e6fb/output.json}{output} \\
\nolinkurl{colette-le-ble-en-herbe-aa5c7bfe3b6367b5} & fr & 4/1/2/0 & reported & \href{https://github.com/mannyrayner/let_me_count_the_ways/blob/662f77a398e884f9bf30601c3bbf433a722c9619/results/annotation/canonical_31_v0_11_v0_3_1/annotations/colette-le-ble-en-herbe-aa5c7bfe3b6367b5/ad7d78d014d68c8c3c4b514a19dbb26672153fa8ea744200a1fa508e2ea19727/request.json}{input}; \href{https://github.com/mannyrayner/let_me_count_the_ways/blob/662f77a398e884f9bf30601c3bbf433a722c9619/results/annotation/canonical_31_v0_11_v0_3_1/annotations/colette-le-ble-en-herbe-aa5c7bfe3b6367b5/ad7d78d014d68c8c3c4b514a19dbb26672153fa8ea744200a1fa508e2ea19727/output.json}{output} \\
\nolinkurl{lagerlof-gosta-berlings-saga-ea36a229384cb926} & sv & 4/0/2/0 & imagined & \href{https://github.com/mannyrayner/let_me_count_the_ways/blob/662f77a398e884f9bf30601c3bbf433a722c9619/results/annotation/canonical_31_v0_11_v0_3_1/annotations/lagerlof-gosta-berlings-saga-ea36a229384cb926/ffd2799451b89b3952e0352b95a167c2cdff2c7c183dbda925b15794c5697ac8/request.json}{input}; \href{https://github.com/mannyrayner/let_me_count_the_ways/blob/662f77a398e884f9bf30601c3bbf433a722c9619/results/annotation/canonical_31_v0_11_v0_3_1/annotations/lagerlof-gosta-berlings-saga-ea36a229384cb926/ffd2799451b89b3952e0352b95a167c2cdff2c7c183dbda925b15794c5697ac8/output.json}{output} \\
\nolinkurl{strindberg-froken-julie-1492b64f21186bd3} & sv & 3/0/0/0 & quoted or revoiced & \href{https://github.com/mannyrayner/let_me_count_the_ways/blob/662f77a398e884f9bf30601c3bbf433a722c9619/results/annotation/canonical_31_v0_11_v0_3_1/annotations/strindberg-froken-julie-1492b64f21186bd3/69f514c5c7614afaedafaf686242a7e26dec38246171e6e867b26ec1a9a101de/request.json}{input}; \href{https://github.com/mannyrayner/let_me_count_the_ways/blob/662f77a398e884f9bf30601c3bbf433a722c9619/results/annotation/canonical_31_v0_11_v0_3_1/annotations/strindberg-froken-julie-1492b64f21186bd3/69f514c5c7614afaedafaf686242a7e26dec38246171e6e867b26ec1a9a101de/output.json}{output} \\
\nolinkurl{rostand-cyrano-de-bergerac-920ae26958d61661} & fr & 1/0/0/0 & hypothetical & \href{https://github.com/mannyrayner/let_me_count_the_ways/blob/662f77a398e884f9bf30601c3bbf433a722c9619/results/annotation/canonical_31_v0_11_v0_3_1/annotations/rostand-cyrano-de-bergerac-920ae26958d61661/65add50ff86ab6a7de1aab7891554cdd277c166104396addb6143ea84f744bfd/request.json}{input}; \href{https://github.com/mannyrayner/let_me_count_the_ways/blob/662f77a398e884f9bf30601c3bbf433a722c9619/results/annotation/canonical_31_v0_11_v0_3_1/annotations/rostand-cyrano-de-bergerac-920ae26958d61661/65add50ff86ab6a7de1aab7891554cdd277c166104396addb6143ea84f744bfd/output.json}{output} \\
\nolinkurl{fontane-effi-briest-c9a7185a57602a96} & de & 4/0/0/0 & direct & \href{https://github.com/mannyrayner/let_me_count_the_ways/blob/662f77a398e884f9bf30601c3bbf433a722c9619/results/annotation/canonical_31_v0_11_v0_3_1/annotations/fontane-effi-briest-c9a7185a57602a96/ce8f372f2ae2f3f263e36ecd4fb3937c10242b8647caa2d349f34a74e3614e78/request.json}{input}; \href{https://github.com/mannyrayner/let_me_count_the_ways/blob/662f77a398e884f9bf30601c3bbf433a722c9619/results/annotation/canonical_31_v0_11_v0_3_1/annotations/fontane-effi-briest-c9a7185a57602a96/ce8f372f2ae2f3f263e36ecd4fb3937c10242b8647caa2d349f34a74e3614e78/output.json}{output} \\
\nolinkurl{bronte-jane-eyre-5da4cfc02f8bc2c8} & en & 4/0/0/0 & direct & \href{https://github.com/mannyrayner/let_me_count_the_ways/blob/662f77a398e884f9bf30601c3bbf433a722c9619/results/annotation/canonical_31_v0_11_v0_3_1/annotations/bronte-jane-eyre-5da4cfc02f8bc2c8/c03369f533e2232b1ded4217e8d6f72de5f1facc2fc907eff489feec9fec6e1a/request.json}{input}; \href{https://github.com/mannyrayner/let_me_count_the_ways/blob/662f77a398e884f9bf30601c3bbf433a722c9619/results/annotation/canonical_31_v0_11_v0_3_1/annotations/bronte-jane-eyre-5da4cfc02f8bc2c8/c03369f533e2232b1ded4217e8d6f72de5f1facc2fc907eff489feec9fec6e1a/output.json}{output} \\
\nolinkurl{rostand-cyrano-de-bergerac-ca838af570a46ee9} & fr & 4/0/0/0 & direct & \href{https://github.com/mannyrayner/let_me_count_the_ways/blob/662f77a398e884f9bf30601c3bbf433a722c9619/results/annotation/canonical_31_v0_11_v0_3_1/annotations/rostand-cyrano-de-bergerac-ca838af570a46ee9/45bda8dba9cf1b2c99d170d9f5cdbf74a51fe9ff9830975fcda37e3e6f74c975/request.json}{input}; \href{https://github.com/mannyrayner/let_me_count_the_ways/blob/662f77a398e884f9bf30601c3bbf433a722c9619/results/annotation/canonical_31_v0_11_v0_3_1/annotations/rostand-cyrano-de-bergerac-ca838af570a46ee9/45bda8dba9cf1b2c99d170d9f5cdbf74a51fe9ff9830975fcda37e3e6f74c975/output.json}{output} \\
\nolinkurl{zuccoli-lamore-di-loredana-9a817d5f2cd7e04d} & it & 4/0/0/0 & direct & \href{https://github.com/mannyrayner/let_me_count_the_ways/blob/662f77a398e884f9bf30601c3bbf433a722c9619/results/annotation/canonical_31_v0_11_v0_3_1/annotations/zuccoli-lamore-di-loredana-9a817d5f2cd7e04d/d147be48638a2c087d8037848fd91b61b8de0f659271ae456768755f49ba2f7e/request.json}{input}; \href{https://github.com/mannyrayner/let_me_count_the_ways/blob/662f77a398e884f9bf30601c3bbf433a722c9619/results/annotation/canonical_31_v0_11_v0_3_1/annotations/zuccoli-lamore-di-loredana-9a817d5f2cd7e04d/d147be48638a2c087d8037848fd91b61b8de0f659271ae456768755f49ba2f7e/output.json}{output} \\
\nolinkurl{ibsen-et-dukkehjem-22aa549e55c520e4} & no & 4/0/0/0 & direct & \href{https://github.com/mannyrayner/let_me_count_the_ways/blob/662f77a398e884f9bf30601c3bbf433a722c9619/results/annotation/canonical_31_v0_11_v0_3_1/annotations/ibsen-et-dukkehjem-22aa549e55c520e4/d6ccdcad59be3fd7f845ef49fa1b1b5cdbd9f46d8e8d001b4dd544b0aa2be3cd/request.json}{input}; \href{https://github.com/mannyrayner/let_me_count_the_ways/blob/662f77a398e884f9bf30601c3bbf433a722c9619/results/annotation/canonical_31_v0_11_v0_3_1/annotations/ibsen-et-dukkehjem-22aa549e55c520e4/d6ccdcad59be3fd7f845ef49fa1b1b5cdbd9f46d8e8d001b4dd544b0aa2be3cd/output.json}{output} \\
\nolinkurl{lagerlof-gosta-berlings-saga-baef76cd44bc6283} & sv & 4/0/0/0 & direct & \href{https://github.com/mannyrayner/let_me_count_the_ways/blob/662f77a398e884f9bf30601c3bbf433a722c9619/results/annotation/canonical_31_v0_11_v0_3_1/annotations/lagerlof-gosta-berlings-saga-baef76cd44bc6283/2f3d27b79060d241ea9bab7bb7db947b61771002c239216cf55b056001f11a77/request.json}{input}; \href{https://github.com/mannyrayner/let_me_count_the_ways/blob/662f77a398e884f9bf30601c3bbf433a722c9619/results/annotation/canonical_31_v0_11_v0_3_1/annotations/lagerlof-gosta-berlings-saga-baef76cd44bc6283/2f3d27b79060d241ea9bab7bb7db947b61771002c239216cf55b056001f11a77/output.json}{output} \\
\nolinkurl{ibsen-et-dukkehjem-003dfd3f3df5ffc6} & no & 4/0/0/0 & direct & \href{https://github.com/mannyrayner/let_me_count_the_ways/blob/662f77a398e884f9bf30601c3bbf433a722c9619/results/annotation/canonical_31_v0_11_v0_3_1/annotations/ibsen-et-dukkehjem-003dfd3f3df5ffc6/be28d317c457e521745f9278ed025a1ba7c36484cee27e43790c82a40991e4b3/request.json}{input}; \href{https://github.com/mannyrayner/let_me_count_the_ways/blob/662f77a398e884f9bf30601c3bbf433a722c9619/results/annotation/canonical_31_v0_11_v0_3_1/annotations/ibsen-et-dukkehjem-003dfd3f3df5ffc6/be28d317c457e521745f9278ed025a1ba7c36484cee27e43790c82a40991e4b3/output.json}{output} \\
\nolinkurl{ibsen-et-dukkehjem-730c4a81efa614e1} & no & 4/0/0/0 & direct & \href{https://github.com/mannyrayner/let_me_count_the_ways/blob/662f77a398e884f9bf30601c3bbf433a722c9619/results/annotation/canonical_31_v0_11_v0_3_1/annotations/ibsen-et-dukkehjem-730c4a81efa614e1/40861f39107147506c3a747e6dbf84c09f4c18a8c76c45bd67865def4ed0f261/request.json}{input}; \href{https://github.com/mannyrayner/let_me_count_the_ways/blob/662f77a398e884f9bf30601c3bbf433a722c9619/results/annotation/canonical_31_v0_11_v0_3_1/annotations/ibsen-et-dukkehjem-730c4a81efa614e1/40861f39107147506c3a747e6dbf84c09f4c18a8c76c45bd67865def4ed0f261/output.json}{output} \\
\nolinkurl{flaubert-madame-bovary-176b95833345dab6} & fr & 4/0/2/0 & direct & \href{https://github.com/mannyrayner/let_me_count_the_ways/blob/662f77a398e884f9bf30601c3bbf433a722c9619/results/annotation/canonical_31_v0_11_v0_3_1/annotations/flaubert-madame-bovary-176b95833345dab6/48b27c1bf065d9e4892070669f33a8a96cb8b29ee63bdc5c8b35a1c8c3e1b952/request.json}{input}; \href{https://github.com/mannyrayner/let_me_count_the_ways/blob/662f77a398e884f9bf30601c3bbf433a722c9619/results/annotation/canonical_31_v0_11_v0_3_1/annotations/flaubert-madame-bovary-176b95833345dab6/48b27c1bf065d9e4892070669f33a8a96cb8b29ee63bdc5c8b35a1c8c3e1b952/output.json}{output} \\
\nolinkurl{flaubert-madame-bovary-c648bd2e9a7ad647} & fr & 4/0/3/0 & quoted or revoiced & \href{https://github.com/mannyrayner/let_me_count_the_ways/blob/662f77a398e884f9bf30601c3bbf433a722c9619/results/annotation/canonical_31_v0_11_v0_3_1/annotations/flaubert-madame-bovary-c648bd2e9a7ad647/fbfa1810e090f58e0094985a122d76b8cda1aa061d117e1a0cddc8476e2bbb99/request.json}{input}; \href{https://github.com/mannyrayner/let_me_count_the_ways/blob/662f77a398e884f9bf30601c3bbf433a722c9619/results/annotation/canonical_31_v0_11_v0_3_1/annotations/flaubert-madame-bovary-c648bd2e9a7ad647/fbfa1810e090f58e0094985a122d76b8cda1aa061d117e1a0cddc8476e2bbb99/output.json}{output} \\
\nolinkurl{flaubert-madame-bovary-d3ce852d259a04be} & fr & 4/0/2/0 & direct & \href{https://github.com/mannyrayner/let_me_count_the_ways/blob/662f77a398e884f9bf30601c3bbf433a722c9619/results/annotation/canonical_31_v0_11_v0_3_1/annotations/flaubert-madame-bovary-d3ce852d259a04be/8f8424ca127ae4b42f93013197295648eed6d1a221e80a5abf8538bd737e493e/request.json}{input}; \href{https://github.com/mannyrayner/let_me_count_the_ways/blob/662f77a398e884f9bf30601c3bbf433a722c9619/results/annotation/canonical_31_v0_11_v0_3_1/annotations/flaubert-madame-bovary-d3ce852d259a04be/8f8424ca127ae4b42f93013197295648eed6d1a221e80a5abf8538bd737e493e/output.json}{output} \\
\nolinkurl{nansen-maria-93ea665425fd352f} & da & 4/4/0/0 & direct & \href{https://github.com/mannyrayner/let_me_count_the_ways/blob/662f77a398e884f9bf30601c3bbf433a722c9619/results/annotation/canonical_31_v0_11_v0_3_1/annotations/nansen-maria-93ea665425fd352f/d5338f7ac0dbdd0b998461aa6ff8cd9fbc280d1a43db67ca5e57088f6d7d8c75/request.json}{input}; \href{https://github.com/mannyrayner/let_me_count_the_ways/blob/662f77a398e884f9bf30601c3bbf433a722c9619/results/annotation/canonical_31_v0_11_v0_3_1/annotations/nansen-maria-93ea665425fd352f/d5338f7ac0dbdd0b998461aa6ff8cd9fbc280d1a43db67ca5e57088f6d7d8c75/output.json}{output} \\
\nolinkurl{nansen-maria-a95aa1804441024d} & da & 4/4/0/0 & direct & \href{https://github.com/mannyrayner/let_me_count_the_ways/blob/662f77a398e884f9bf30601c3bbf433a722c9619/results/annotation/canonical_31_v0_11_v0_3_1/annotations/nansen-maria-a95aa1804441024d/6dbeb2f2a84e2012c634ccc046d6adc3bc5a745d7ade75aed3dd68ee5778a2b0/request.json}{input}; \href{https://github.com/mannyrayner/let_me_count_the_ways/blob/662f77a398e884f9bf30601c3bbf433a722c9619/results/annotation/canonical_31_v0_11_v0_3_1/annotations/nansen-maria-a95aa1804441024d/6dbeb2f2a84e2012c634ccc046d6adc3bc5a745d7ade75aed3dd68ee5778a2b0/output.json}{output} \\
\nolinkurl{bronte-tenant-of-wildfell-hall-14275eeb09f3d72c} & en & 3/4/0/0 & direct & \href{https://github.com/mannyrayner/let_me_count_the_ways/blob/662f77a398e884f9bf30601c3bbf433a722c9619/results/annotation/canonical_31_v0_11_v0_3_1/annotations/bronte-tenant-of-wildfell-hall-14275eeb09f3d72c/0726c9eb66614dc79758da6cdfe17d0bf07ad9d1313e53607ff7caa70a803934/request.json}{input}; \href{https://github.com/mannyrayner/let_me_count_the_ways/blob/662f77a398e884f9bf30601c3bbf433a722c9619/results/annotation/canonical_31_v0_11_v0_3_1/annotations/bronte-tenant-of-wildfell-hall-14275eeb09f3d72c/0726c9eb66614dc79758da6cdfe17d0bf07ad9d1313e53607ff7caa70a803934/output.json}{output} \\
\end{longtable}
\endgroup

\fi
\end{document}
